\chapter{Traits \code{GraphRead} / \code{GraphWrite}}
\label{ch:graph-traits}

\begin{note}
Stub. El refactor P2-B extrajo estos traits para que los
consumidores del motor (scorer de anomalía, extractor de
features GNN, writer de ingesta) puedan depender de una API
semántica en vez del struct concreto \code{GraphHunter}. El
objetivo eventual es poder enchufar backends remotos (Neo4j, un
store MVCC streaming) sin editar el scorer, el matcher ni el
writer --- con tal de que existan las impls del trait.
\end{note}

\section{Traits}

\begin{lstlisting}[style=rust]
pub trait GraphRead {
    fn entity_count(&self) -> usize;
    fn relation_count(&self) -> usize;
    fn get_entity(&self, id: &str) -> Option<&Entity>;
    fn get_compact_relations(&self, src: &str)
        -> Cow<'_, [CompactRelation]>;
    fn get_reverse_source_ids(&self, id: &str) -> Vec<String>;
    fn entity_ids_for_type(&self, et: &EntityType) -> Option<Vec<String>>;
    fn entity_types_in_graph(&self) -> Vec<String>;
    fn materialize_relation(&self, c: &CompactRelation) -> Relation;
    fn supports_simd(&self) -> bool { false }
}

pub trait GraphWrite: GraphRead {
    fn add_entity(&mut self, e: Entity)   -> Result<(), GraphError>;
    fn add_relation(&mut self, r: Relation) -> Result<(), GraphError>;
}
\end{lstlisting}

El retorno \code{Cow<'\_, [CompactRelation]>} es la clave que
permite coexistencia de backends remotos:
\code{InMemoryGraph} devuelve \code{Cow::Borrowed(\&slice)} con
costo cero; una impl remota devuelve \code{Cow::Owned(vec)}
cuando tiene que sintetizar la lista.

\section{Qué va en el trait y qué no}

\textbf{Va:} lectura y escritura semánticas, las operaciones que
un hipotético backend remoto necesitaría reimplementar.

\textbf{No va:} el matcher en sí (\code{search\_temporal\_pattern}),
el despacho SIMD, los analíticos (\code{pagerank},
\code{betweenness}). Son específicos del motor; un backend Neo4j
traería su propio motor de patrones en vez de implementar el DFS
de Rust contra un \code{\&dyn GraphRead}.

\section{Contrato dual-impl}

\filepath{graph\_hunter\_core/tests/backend\_dual\_impl.rs} corre
un arnés genérico \code{validate\_backend<B: GraphWrite>(b)} que
ejercita cada método del trait contra \code{GraphHunter}. Cuando
aterrice una segunda impl, el mismo arnés la cubre gratis.

\begin{inthecode}
\filepath{graph\_hunter\_core/src/backend/mod.rs} --- definiciones
del trait.\\
\filepath{graph\_hunter\_core/src/backend/in\_memory.rs} ---
\code{impl GraphRead + GraphWrite for GraphHunter}.\\
\filepath{graph\_hunter\_core/tests/backend\_dual\_impl.rs} ---
validación dual-impl.
\end{inthecode}
